% ===========================================================================
% Category: Density-Based Learning
% Generated from ML_Lab_Manual_Completed.md - edit the Markdown and re-run
% scripts/build_latex_manual.py instead of editing this file by hand.
% ===========================================================================

\chapter{Experiment 5: DBSCAN}\label{experiment-5-dbscan}

\textbf{Category:} Density-based learning

\section{1. Objective}\label{objective}

Detect dense groups and noise with DBSCAN without specifying the number
of clusters, on real spatial data.

\section{2. Dataset Used}\label{dataset-used}

Real USGS earthquake catalogue 2023
(\passthrough{\lstinline!data/earthquakes.csv!}): 16,190 global events
with magnitude ≥ 4.0; the experiment uses the 7,638 events with
magnitude ≥ 4.5. Coordinates are converted to radians and clustered with
the haversine (great-circle) metric.

\section{3. Required Libraries}\label{required-libraries}

scikit-learn, numpy, pandas, matplotlib

\section{4. Brief Theory}\label{brief-theory}

DBSCAN uses a radius ε and minimum points MinPts. Points with ≥ MinPts
neighbors within ε are core points; points within ε of a core are border
points; the rest are noise (label −1). Clusters are chains of
density-reachable core points, so arbitrary shapes can be discovered. A
single global ε struggles when densities vary.

\section{5. Procedure}\label{procedure}

\begin{enumerate}
\def\labelenumi{\arabic{enumi}.}
\tightlist
\item
  Load the catalogue and filter to magnitude ≥ 4.5 (well-recorded
  events).
\item
  Convert latitude/longitude to radians for the haversine metric.
\item
  Choose ε from the k-distance graph (k = 10) --- the knee at 0.03 rad ≈
  191 km.
\item
  Fit DBSCAN and count clusters/noise.
\item
  Show a parameter-sensitivity table for at least two settings.
\item
  Visualize epicenters and check the largest clusters against known
  seismic regions.
\end{enumerate}

\section{6. Reference Implementation}\label{reference-implementation}

\begin{lstlisting}[language=Python]
X_rad = np.radians(quakes[['latitude', 'longitude']].values)
nbrs = NearestNeighbors(n_neighbors=10, metric='haversine').fit(X_rad)
k_distances = np.sort(nbrs.kneighbors(X_rad)[0][:, -1])       # knee ~0.03 rad
dbscan = DBSCAN(eps=0.03, min_samples=10, metric='haversine')
db_labels = dbscan.fit_predict(X_rad)
\end{lstlisting}

\section{7. Results}\label{results}

Parameter sensitivity (min\_samples = 10):

{\def\LTcaptype{none} % do not increment counter
\begin{longtable}[]{@{}llll@{}}
\toprule\noalign{}
eps (rad) & approx km & clusters & noise \% \\
\midrule\noalign{}
\endhead
\bottomrule\noalign{}
\endlastfoot
0.02 & 127 & 70 & 21.2\% \\
\textbf{0.03 (selected)} & \textbf{191} & \textbf{55} &
\textbf{13.5\%} \\
0.05 & 319 & 50 & 5.7\% \\
\end{longtable}
}

Largest clusters (real-world validation): Philippines (1,629 events),
Papua New Guinea (1,068), Tonga (947), Japan (895), South Sandwich
Islands (253), Turkey (207). Silhouette on non-noise (3,000-event
sample) = 0.326.

\section{8. Discussion and Conclusion}\label{discussion-and-conclusion}

DBSCAN recovers the global seismic structure without labels: dense
chains along tectonic plate boundaries form clusters, while
\textasciitilde1,000 isolated intraplate events are noise. Smaller ε
fragments clusters and labels more events as noise; larger ε merges
boundaries and reduces noise. The single global ε is the main
limitation, which motivates HDBSCAN.

\section{9. Answers to Viva Questions}\label{answers-to-viva-questions}

\begin{itemize}
\tightlist
\item
  \textbf{Why does DBSCAN not require K?} Clusters are grown from
  density-connected core points, so their number and shape emerge from
  the data instead of being fixed in advance.
\item
  \textbf{What does label −1 mean?} The point is neither a core nor a
  border point --- it lies in a sparse region and is treated as
  noise/outlier (here: isolated earthquakes).
\item
  \textbf{Why can DBSCAN struggle with varying densities?} A single ε
  that is small enough for a dense cluster fragments a sparse one, while
  one large enough for the sparse cluster merges the dense ones --- no
  global setting fits both.
\end{itemize}

\begin{center}\rule{0.5\linewidth}{0.5pt}\end{center}

\section{10. Complete Code Listing}\label{complete-code-listing}

\emph{Full runnable program for this experiment, extracted from
\passthrough{\lstinline!notebooks/02\_density\_based\_learning.ipynb!}
(executed; results above are its actual output).}

\textbf{Listing 5.1 --- Core numerical and plotting libraries}

\begin{lstlisting}[language=Python]
# ---- Core numerical and plotting libraries ----
import numpy as np
import pandas as pd
import matplotlib.pyplot as plt
import seaborn as sns

# ---- Density-based clustering algorithms ----
from sklearn.cluster import DBSCAN, HDBSCAN
# ---- Nearest-neighbor search (used to pick epsilon) ----
from sklearn.neighbors import NearestNeighbors
# ---- Evaluation metrics ----
from sklearn.metrics import silhouette_score

# Aesthetics setup (consistent look across all notebooks)
plt.style.use('seaborn-v0_8-whitegrid' if 'seaborn-v0_8-whitegrid' in plt.style.available else 'default')
plt.rcParams['font.sans-serif'] = 'DejaVu Sans'
plt.rcParams['figure.figsize'] = (11, 6)
plt.rcParams['figure.dpi'] = 110

np.random.seed(42)
print("Density-based libraries loaded successfully!")
\end{lstlisting}

\textbf{Listing 5.2 --- Load the real USGS earthquake catalogue (global,
2023)}

\begin{lstlisting}[language=Python]
# ---- Load the real USGS earthquake catalogue (global, 2023) ----
df = pd.read_csv('../data/earthquakes.csv')
print(f"Full catalogue (magnitude >= 4.0): {len(df):,} events")
print(df.head(3).to_string(index=False))

# Keep well-recorded moderate+ events: magnitude >= 4.5
quakes = df[df['mag'] >= 4.5].reset_index(drop=True)
print(f"\nAnalysed subset (magnitude >= 4.5): {len(quakes):,} events")

# ---- Convert epicenters to radians [latitude, longitude] ----
# scikit-learn's 'haversine' metric expects radian coordinates and computes
# great-circle distance on the sphere (correct for global data).
X_rad = np.radians(quakes[['latitude', 'longitude']].values)
print("Coordinate range (deg): lat", quakes.latitude.min(), "to", quakes.latitude.max(),
      "| lon", quakes.longitude.min(), "to", quakes.longitude.max())
\end{lstlisting}

\textbf{Listing 5.3 --- Choosing epsilon from the k-distance graph
(haversine metric)}

\begin{lstlisting}[language=Python]
# ---- Choosing epsilon from the k-distance graph (haversine metric) ----
# Sorted distance to each event's k-th nearest neighbor. The knee marks the
# transition from dense seismic zones to sparse, isolated events.
min_samples = 10
nbrs = NearestNeighbors(n_neighbors=min_samples, metric='haversine').fit(X_rad)
distances, _ = nbrs.kneighbors(X_rad)
k_distances = np.sort(distances[:, -1])   # radians (0.03 rad = 191 km on Earth)

plt.figure(figsize=(9, 5))
plt.plot(k_distances, color='crimson', linewidth=2.5)
plt.title(f"k-Distance Graph (k={min_samples}) for Epsilon Selection", fontsize=14, fontweight='bold')
plt.xlabel("Earthquakes Sorted by Distance", fontsize=12)
plt.ylabel(f"{min_samples}-NN Distance (radians)", fontsize=12)
plt.axhline(y=0.03, color='navy', linestyle='--', label='Selected ε = 0.03 rad (~191 km)')
plt.legend()
plt.tight_layout()
plt.show()
\end{lstlisting}

\textbf{Listing 5.4 --- Fit DBSCAN with the chosen epsilon and MinPts}

\begin{lstlisting}[language=Python]
# ---- Fit DBSCAN with the chosen epsilon and MinPts ----
# DBSCAN grows clusters from core events, connects them via density-reachability,
# and labels everything else as noise (-1).
eps_selected = 0.03   # radians -> 0.03 * 6371 km ~ 191 km
dbscan = DBSCAN(eps=eps_selected, min_samples=min_samples, metric='haversine')
db_labels = dbscan.fit_predict(X_rad)

# Identify which events are core points (returned by the model) vs. border points
core_samples_mask = np.zeros_like(db_labels, dtype=bool)
core_samples_mask[dbscan.core_sample_indices_] = True

n_clusters_db = len(set(db_labels)) - (1 if -1 in db_labels else 0)
n_noise_db = np.sum(db_labels == -1)
print(f"DBSCAN: {n_clusters_db} clusters, {n_noise_db:,} noise events ({n_noise_db / len(X_rad) * 100:.1f}%)")

# ---- Visualize epicenters: clustered events in color, noise as black crosses ----
plt.figure(figsize=(12, 5.5))
palette = sns.color_palette("tab20", max(n_clusters_db, 1))
for k in range(n_clusters_db):
    mask = db_labels == k
    plt.scatter(quakes.longitude[mask], quakes.latitude[mask], s=6, color=palette[k % 20], alpha=0.75)

noise_mask = db_labels == -1
plt.scatter(quakes.longitude[noise_mask], quakes.latitude[noise_mask], s=6, color='black', marker='x', label='Noise (-1)')
plt.title(f"DBSCAN on Global Earthquakes 2023 (ε={eps_selected} rad, {n_clusters_db} clusters)", fontsize=13, fontweight='bold')
plt.xlabel("Longitude", fontsize=12)
plt.ylabel("Latitude", fontsize=12)
plt.legend(loc='lower left')
plt.tight_layout()
plt.show()
\end{lstlisting}

\textbf{Listing 5.5 --- Parameter sensitivity: DBSCAN under several
(eps, min\_samples) settings}

\begin{lstlisting}[language=Python]
# ---- Parameter sensitivity: DBSCAN under several (eps, min_samples) settings ----
# The manual requires at least two settings so the effect of eps can be discussed.
print("eps (rad)   approx km   clusters   noise %")
for eps_test in [0.02, 0.03, 0.05]:
    lbl = DBSCAN(eps=eps_test, min_samples=10, metric='haversine').fit_predict(X_rad)
    n_clusters_test = len(set(lbl)) - (1 if -1 in lbl else 0)
    noise_pct = (lbl == -1).mean() * 100
    print(f"  {eps_test:.2f}       {eps_test * 6371:5.0f}      {n_clusters_test:3d}      {noise_pct:5.1f}%")
\end{lstlisting}

\chapter{Experiment 6: HDBSCAN}\label{experiment-6-hdbscan}

\textbf{Category:} Density-based learning

\section{1. Objective}\label{objective-1}

Apply hierarchical density-based clustering to the same spatial data to
discover clusters of different densities and identify noise without
selecting ε.

\section{2. Dataset Used}\label{dataset-used-1}

USGS earthquakes 2023, magnitude ≥ 4.5 (7,638 events), radians +
haversine metric.

\section{3. Required Libraries}\label{required-libraries-1}

scikit-learn (\passthrough{\lstinline!HDBSCAN!}), numpy, pandas,
matplotlib

\section{4. Brief Theory}\label{brief-theory-1}

HDBSCAN builds a hierarchy over all density levels: core distance →
mutual reachability distance d\_mreach(a,b) = max\{d\_core(a),
d\_core(b), d(a,b)\} → minimum spanning tree → condensed cluster tree,
then keeps the most stable clusters S(C) = Σ (λ\_p − λ\_birth(C)). It
also returns membership probabilities, so no global ε is needed.

\section{5. Procedure}\label{procedure-1}

\begin{enumerate}
\def\labelenumi{\arabic{enumi}.}
\tightlist
\item
  Reuse the prepared radian coordinates.
\item
  Fit HDBSCAN with \passthrough{\lstinline!min\_cluster\_size=50!},
  \passthrough{\lstinline!min\_samples=10!},
  \passthrough{\lstinline!metric='haversine'!}.
\item
  Count clusters and noise; inspect membership probabilities.
\item
  Visualize the partition.
\item
  Compare with DBSCAN.
\item
  Discuss the effect of \passthrough{\lstinline!min\_cluster\_size!}.
\end{enumerate}

\section{6. Reference Implementation}\label{reference-implementation-1}

\begin{lstlisting}[language=Python]
hdb = HDBSCAN(min_cluster_size=50, min_samples=10, metric='haversine', copy=False)
hdb_labels = hdb.fit_predict(X_rad)
\end{lstlisting}

\section{7. Results}\label{results-1}

{\def\LTcaptype{none} % do not increment counter
\begin{longtable}[]{@{}
  >{\raggedright\arraybackslash}p{(\linewidth - 8\tabcolsep) * \real{0.2000}}
  >{\raggedright\arraybackslash}p{(\linewidth - 8\tabcolsep) * \real{0.2000}}
  >{\raggedright\arraybackslash}p{(\linewidth - 8\tabcolsep) * \real{0.2000}}
  >{\raggedright\arraybackslash}p{(\linewidth - 8\tabcolsep) * \real{0.2000}}
  >{\raggedright\arraybackslash}p{(\linewidth - 8\tabcolsep) * \real{0.2000}}@{}}
\toprule\noalign{}
\begin{minipage}[b]{\linewidth}\raggedright
Algorithm
\end{minipage} & \begin{minipage}[b]{\linewidth}\raggedright
Clusters
\end{minipage} & \begin{minipage}[b]{\linewidth}\raggedright
Noise events
\end{minipage} & \begin{minipage}[b]{\linewidth}\raggedright
Largest cluster
\end{minipage} & \begin{minipage}[b]{\linewidth}\raggedright
Silhouette (3k sample)
\end{minipage} \\
\midrule\noalign{}
\endhead
\bottomrule\noalign{}
\endlastfoot
DBSCAN (ε=0.03) & 55 & 13.5\% & 1,629 & 0.326 \\
HDBSCAN & 40 & 21.5\% & 560 & \textbf{0.631} \\
\end{longtable}
}

Membership probabilities are high (\textgreater{} 0.9) for events deep
inside seismic zones and lower near cluster edges. HDBSCAN's silhouette
is roughly double DBSCAN's because adaptive density splits DBSCAN's long
arc-shaped clusters (e.g., the whole Japan-Kuril chain) into compact,
well-separated groups.

\section{8. Discussion and
Conclusion}\label{discussion-and-conclusion-1}

HDBSCAN needs no ε and handles the varying density along plate
boundaries better than DBSCAN; it is more conservative (21.5\% vs 13.5\%
noise) because clusters must be stable across density scales. Increasing
\passthrough{\lstinline!min\_cluster\_size!} reduces the number of small
clusters; decreasing it fragments. The cost is a more complex algorithm
and less transparent parameter meaning.

\section{9. Answers to Viva
Questions}\label{answers-to-viva-questions-1}

\begin{itemize}
\tightlist
\item
  \textbf{What is hierarchical about HDBSCAN?} It constructs a full
  hierarchy of density-based clusters across all ε values, then
  condenses it by tracking cluster birth/death and keeping the most
  stable clusters.
\item
  \textbf{What does membership probability mean?} A per-point confidence
  that it belongs to its assigned cluster, derived from the point's
  persistence in the cluster as density thresholds change (low values
  indicate borderline points).
\item
  \textbf{Why is HDBSCAN useful for variable-density data?} Because
  clusters are selected by stability across density levels rather than
  by one global ε, so dense and sparse clusters can coexist.
\end{itemize}

\begin{center}\rule{0.5\linewidth}{0.5pt}\end{center}

\section{10. Complete Code Listing}\label{complete-code-listing-1}

\emph{Full runnable program for this experiment, extracted from
\passthrough{\lstinline!notebooks/02\_density\_based\_learning.ipynb!}
(executed; results above are its actual output).}

\textbf{Listing 6.1 --- Core numerical and plotting libraries}

\begin{lstlisting}[language=Python]
# ---- Core numerical and plotting libraries ----
import numpy as np
import pandas as pd
import matplotlib.pyplot as plt
import seaborn as sns

# ---- Density-based clustering algorithms ----
from sklearn.cluster import DBSCAN, HDBSCAN
# ---- Nearest-neighbor search (used to pick epsilon) ----
from sklearn.neighbors import NearestNeighbors
# ---- Evaluation metrics ----
from sklearn.metrics import silhouette_score

# Aesthetics setup (consistent look across all notebooks)
plt.style.use('seaborn-v0_8-whitegrid' if 'seaborn-v0_8-whitegrid' in plt.style.available else 'default')
plt.rcParams['font.sans-serif'] = 'DejaVu Sans'
plt.rcParams['figure.figsize'] = (11, 6)
plt.rcParams['figure.dpi'] = 110

np.random.seed(42)
print("Density-based libraries loaded successfully!")
\end{lstlisting}

\textbf{Listing 6.2 --- Load the real USGS earthquake catalogue (global,
2023)}

\begin{lstlisting}[language=Python]
# ---- Load the real USGS earthquake catalogue (global, 2023) ----
df = pd.read_csv('../data/earthquakes.csv')
print(f"Full catalogue (magnitude >= 4.0): {len(df):,} events")
print(df.head(3).to_string(index=False))

# Keep well-recorded moderate+ events: magnitude >= 4.5
quakes = df[df['mag'] >= 4.5].reset_index(drop=True)
print(f"\nAnalysed subset (magnitude >= 4.5): {len(quakes):,} events")

# ---- Convert epicenters to radians [latitude, longitude] ----
# scikit-learn's 'haversine' metric expects radian coordinates and computes
# great-circle distance on the sphere (correct for global data).
X_rad = np.radians(quakes[['latitude', 'longitude']].values)
print("Coordinate range (deg): lat", quakes.latitude.min(), "to", quakes.latitude.max(),
      "| lon", quakes.longitude.min(), "to", quakes.longitude.max())
\end{lstlisting}

\textbf{Listing 6.3 --- Fit HDBSCAN (no global epsilon required)}

\begin{lstlisting}[language=Python]
# ---- Fit HDBSCAN (no global epsilon required) ----
# min_cluster_size = smallest group considered a cluster; min_samples controls conservativeness.
hdb = HDBSCAN(min_cluster_size=50, min_samples=10, metric='haversine', copy=False)
hdb_labels = hdb.fit_predict(X_rad)

n_clusters_hdb = len(set(hdb_labels)) - (1 if -1 in hdb_labels else 0)
n_noise_hdb = np.sum(hdb_labels == -1)
print(f"HDBSCAN: {n_clusters_hdb} clusters, {n_noise_hdb:,} noise events ({n_noise_hdb / len(X_rad) * 100:.1f}%)")

# ---- Visualize the HDBSCAN partition ----
plt.figure(figsize=(12, 5.5))
palette = sns.color_palette("tab20", max(n_clusters_hdb, 1))
for k in range(n_clusters_hdb):
    mask = hdb_labels == k
    plt.scatter(quakes.longitude[mask], quakes.latitude[mask], s=6, color=palette[k % 20], alpha=0.75)

noise_mask = hdb_labels == -1
plt.scatter(quakes.longitude[noise_mask], quakes.latitude[noise_mask], s=6, color='black', marker='x', label='Noise (-1)')
plt.title(f"HDBSCAN on Global Earthquakes 2023 ({n_clusters_hdb} clusters, no ε)", fontsize=13, fontweight='bold')
plt.xlabel("Longitude", fontsize=12)
plt.ylabel("Latitude", fontsize=12)
plt.legend(loc='lower left')
plt.tight_layout()
plt.show()
\end{lstlisting}
